\begin{description}
\item[a.] Assume that the average Earth's orbit is at the Sun. Therefore, the
  average great conjunction period equals to the synodic period between Jupiter
  and Saturn.
  \begin{align*}
  \frac{1}{P_{\mathrm{Synodic}}} &= \frac{1}{P_{\mathrm{Jupiter}}} - \frac{1}{P_{\mathrm{Saturn}}} \tag*{[3.0]}\\
  \frac{1}{P_{\mathrm{Synodic}}} &= \frac{1}{11.86} - \frac{1}{29.45}\\
  P_{\mathrm{Synodic}} &= 19.86~\text{years} \tag*{[1.0]}
  \end{align*}
  From average great conjunction period, in 19.86 years Jupiter and Saturn move
  \[
  \text{Heliocentric angle} = \frac{19.86}{11.86}\times360^\circ
    = 602.7^\circ = 242.7^\circ
  \]
  in an eastward direction through the zodiac, or
  $360^\circ - 242.7^\circ = 117.3^\circ$ in a westward direction through the
  zodiac. \hfill[2.0]
\item[b.] On 21st December 2020 (winter solstice), the Sun will be in the
  constellation of Sagittarius. Planets are always near the ecliptic. The
  current 12 zodiac constellations can be roughly divided into 12 equal parts
  of 30 degrees each.

  The conjunction will locate elongation $30.3^\circ$ East to the Sun.
  % sic: "will locate elongation" as in the original
  Therefore, the conjunction will be in the constellation of
  \textbf{Capricornus (Cap)}. \hfill[2.0]
\item[c.] The duration between 7BC and 2020 is 2026 years or 2027 years.

  Each conjunction occurs on the average every 19.86 years. Therefore, the
  number of conjunction between 7BC and 2020 is
  \[
  \text{Integer value of } 2026/19.86 = 102 \tag*{[2.0]}
  \]
  From a conjunction to its previous successive conjunction, the conjunction
  moves $243^\circ$ westwards. The position of the conjunction in 7BC is
  \[
  102\times242.7^\circ = 24758.33^\circ = 278.3^\circ \tag*{[3.0]}
  \]
  % sic: 102 x 242.7 = 24755.4, source prints 24758.33
  westward, or equivalently $81.7^\circ$ eastward from the position of the
  conjunction in 2020. Therefore, the conjunction was in the constellation of
  \textbf{Aries (Ari)} or \textbf{Pisces (Psc)}. \hfill[3.0]
\item[d.] At the second conjunction, Jupiter and Saturn had retrograde
  motions. For observers on the Earth, both of them were in opposition.
  Therefore, the sun was in the constellation of \textbf{Libra (Lib)} (if
  answer c) is Aries) or \textbf{Virgo (Vir)} (if answer c) is Pisces).
  \hfill[4.0]
\end{description}
